\documentclass[11pt,a4paper]{article}
\usepackage[T1]{fontenc}
\usepackage[utf8]{inputenc}
\usepackage{amsmath,amssymb,amsthm}
\usepackage[margin=1in]{geometry}
\usepackage[colorlinks=true,linkcolor=blue,urlcolor=blue]{hyperref}
\theoremstyle{plain}
\newtheorem{theorem}{Theorem}
\newtheorem{lemma}[theorem]{Lemma}
\newtheorem{proposition}[theorem]{Proposition}
\newtheorem{corollary}[theorem]{Corollary}
\theoremstyle{definition}
\newtheorem{remark}[theorem]{Remark}

\title{Recovering the unwritten recurrences of the table OEIS A251616}
\author{Adrian Perez Fontelles\\ \small Independent researcher}
\date{2 September 2026}
\begin{document}
\maketitle

\begin{abstract}
OEIS A251616 is a two-dimensional table $T(n,k)$ whose empirical recurrences are, for several of
its lines, given only as an ORDER --- ``k=2: [order 49]'' and the like --- with the
coefficients in a linked file rather than in the entry. Those conjectures can still be settled,
and the recurrences recovered. A line of the table is a fixed-width or fixed-height array
count, hence a walk count on finitely many vertices, hence satisfies some monic recurrence of
order at most that number; Berlekamp--Massey on twice as many exact terms returns the MINIMAL
one. Where its order equals the order the entry states --- and where the same holds after the
first few terms are dropped --- any recurrence of that order the line satisfies has a
characteristic polynomial that is a multiple of the minimal one and of equal degree, hence is
that polynomial. This note settles column 2.
\end{abstract}

\noindent\small 2020 Mathematics Subject Classification. 05A15, 11B37, 15A18.\normalsize

\section{The table and the unwritten conjectures}

OEIS A251616 is ``T(n,k)=Number of (n+1)X(k+1) 0..6 arrays with every 2X2 subblock summing to a prime.''. It is read by antidiagonals and begins
\[
826, 14008, 14008, 237566, 583034, 237566, 4031056, 24317054,\ \dots
\]
The entry carries, among its empirical claims:
\begin{quote}\small
k=2: [order 49]
\end{quote}
As of the ``Last modified'' line on the live entry (Jul 23 2025 13:28:38, revision 6) these are still
recorded as empirical, and the recurrences themselves are not in the entry text.

\section{Why a line of the table is a walk count}

Fix the index that the line fixes. The entry defines $T(n,k)$ as a number of arrays of a shape
determined by $n$ and $k$, subject to a condition local to a bounded window of consecutive
lines. Substituting the fixed index into the entry's own wording turns the definition into an
ordinary one-parameter array count, and for such a count the admissible windows are the
vertices of a finite digraph and an array is a walk. So the line is a walk count on $S$
vertices, and by the Cayley--Hamilton theorem it satisfies a monic integer recurrence of order
at most $S$.

\section{Recovering the recurrences}

\begin{lemma}\label{lem:bm}
A sequence known to satisfy some linear recurrence of order at most $S$ is determined, with its
minimal recurrence, by its first $2S$ terms; Berlekamp--Massey applied to them returns that
minimal recurrence exactly.
\end{lemma}

\subsection*{Column $k=2$}
Substituting the fixed index into the entry's wording gives an ordinary one-parameter array
count, a walk on $S=343$ vertices. Berlekamp--Massey on $2S=686$ exact terms
returns a minimal recurrence of order $49$ --- the order the entry states --- with
integer coefficients
\[
(c_1,\dots,c_{49})=(61,\ 2656,\ -216061,\ -1387330,\ 280701862,\ -1257360157,\ -188408927989,\ 1887730050237,\ 73784606748362,\ -1040441316624966,\ -17590975813624055,\ 325194870904798787,\ 2514506146011920374,\ -64374099723210953518,\ -186035369688519491044,\ 8484687238494645519719,\ -332076787670414826261,\ -763606912279457261984135,\ 1528891178835577200107819,\ 47369765992817229466574374,\ -163753189358254388221865916,\ -2018099047187261369654612100,\ 9445850842555674961217448440,\ 58053623136081071763010615688,\ -343613069868783372637241984000,\ -1088361018803076224269643294624,\ 8222038484909253437358316466176,\ 12280985271895124459846522240256,\ -131083356288962311583670091765632,\ -61850959315271519862421160666112,\ 1391818660446107313208307007716864,\ -262922385193035991227200050152448,\ -9715085430405446746133694595129344,\ 6189848263376145429582790865965056,\ 43244982729050827141463604964245504,\ -41327499101903474505252770718351360,\ -115882135860670943349543334921175040,\ 138521620979043651078298338556772352,\ 168686045492783207310642718414209024,\ -234965456231140317505992967889879040,\ -113760648183940010779699566178467840,\ 177672274516908918640417426322227200,\ 40254493793291410770415999726387200,\ -57670606862901396995835634384896000,\ -9113798017661595641537164738560000,\ 6526330568291631850158292992000000,\ 880798167809967577475579904000000,\ -100829925636755527670169600000000,\ -12987453604044286092902400000000).
\]
so that $a(m)=\sum_{i=1}^{49}c_i\,a(m-i)$ for every $m$. Removing the first
$49+4$ terms and repeating gives order $49$ again, which is what the
identification below needs.

\begin{theorem}
For each line above, the displayed recurrence holds for every index, and it is the recurrence
the entry records.
\end{theorem}

\begin{proof}
It holds by Lemma~\ref{lem:bm}. For the identification, the entry asserts that the line
satisfies SOME recurrence of the stated order, possibly only from some index on. Let $q$ be its
characteristic polynomial and $q_{\min}$ the minimal polynomial of the tail. Then
$q_{\min}\mid q$, and the second Berlekamp--Massey run --- on the line with its first few
terms removed --- shows $\deg q_{\min}$ equals the stated order, which is $\deg q$. Both are
monic, so $q=q_{\min}$.
\end{proof}

\section{Verification}

Three checks.

First, each line's model was compared against the entry's own published values for that line,
taken from the antidiagonal DATA read in either orientation and from the ``Table starts''
block, and agrees at every available term. This is the only point at which reading the entry's
English enters, and a misreading gives different counts.

Second, each recovered recurrence was evaluated directly on the model's values, with no
Berlekamp--Massey involved, and holds at every index tested.

Third, each order was computed twice by independent routes --- modulo a $61$-bit prime, where
every intermediate is a machine word, and again in exact rational arithmetic --- and once more
on the line with its first few terms dropped, which is what the identification requires. All
agree.

\begin{thebibliography}{9}
\bibitem{oeis} The OEIS Foundation, \emph{The On-Line Encyclopedia of Integer
Sequences}, \url{https://oeis.org}, 2026. Sequence A251616.
\bibitem{massey} J.~L.~Massey, Shift-register synthesis and BCH decoding, \emph{IEEE
Trans. Inform. Theory} \textbf{15} (1969), 122--127.
\bibitem{stanley} R.~P.~Stanley, \emph{Enumerative Combinatorics}, Volume 1, 2nd ed.,
Cambridge University Press, 2011. (Transfer-matrix method, Section 4.7.)
\bibitem{horn} R.~A.~Horn and C.~R.~Johnson, \emph{Matrix Analysis}, 2nd ed., Cambridge
University Press, 2013.
\end{thebibliography}

\end{document}
